\section{BERT: Encoder-Only y Bidirectional Representation}

BERT muestra otra vía: usa solo bloques encoder, permite que cada token lea el contexto en ambas direcciones y aprende una representation mediante masked language modeling. Su interfaz principal no es la generación autorregresiva abierta; en cambio, ofrece contextual embeddings a los que se les puede aplicar fine-tuning para tareas posteriores como clasificación, span prediction, recuperación y token labeling.

\subsection{Masked language modeling}

El \term{masked language modeling} (MLM, modelado de lenguaje enmascarado) elige al azar algunas posiciones de la entrada y le pide al modelo que prediga el token original. Como las posiciones objetivo están enmascaradas, el encoder puede usar el contexto de ambos lados sin leer directamente la respuesta. El BERT original también usaba next sentence prediction (NSP). Los modelos posteriores resolvieron de distintas maneras si NSP era necesaria, pero lo que se ha mantenido como núcleo es la representation bidireccional del encoder.

\lecturefigure{slide-17-bert.jpg}{BERT se apoya en una representación bidireccional encoder-only y en MLM para las tareas posteriores de comprensión.}{Video oficial 20:30; Devlin et al. 2018.}

\begin{knowledgebox}{Lectura de la figura: la diferencia entre BERT y GPT es, ante todo, de visibilidad}
En BERT, cada token puede leer a la vez el contexto izquierdo y el derecho, por lo que es adecuado para aprender la contextual representation de una oración completa. En GPT, cada token causal solo puede leer el prefijo a su izquierda, lo que concuerda de forma natural con la generación paso a paso. El fine-tuning de BERT suele agregar una task head y actualizar los parámetros; con el prompting de GPT, la tarea puede describirse dentro de una interfaz de generación unificada. Ambos provienen del Transformer, pero difieren en la mask, el objective y la product interface.
\end{knowledgebox}

\begin{importantbox}{Comparación entre encoder-only y decoder-only}
\begin{tabular}{L{0.18\textwidth} L{0.36\textwidth} L{0.36\textwidth}}
\toprule
Dimensión & Encoder-only (BERT) & Decoder-only (GPT) \\
\midrule
Visibilidad & self-attention bidireccional & causal left-to-right attention \\
Objetivo de preentrenamiento & masked token prediction & next-token prediction \\
Salida típica & contextual representation & token distribution / generation \\
Interfaz de la tarea & fine-tuning, linear head, embedding & prompting, completion, fine-tuning \\
Ventaja & tareas intensivas en comprensión y representación & interfaz unificada de generación e interacción \\
\bottomrule
\end{tabular}
\end{importantbox}

\teachervoice{El ponente cierra la introducción con BERT y enfatiza el bloque encoder, la masked prediction y el fine-tuning para tareas posteriores. Este orden es importante: primero se separan dos direcciones a partir del encoder-decoder original y después se usan GPT y BERT para explicar las familias de arquitecturas, en lugar de deducir el mecanismo a partir de los nombres comerciales.}

\subsection{Resumen de la sección}

El núcleo de BERT es la representation bidireccional del encoder y la masked prediction. Comparte con GPT el bloque de self-attention, pero sirve a interfaces de tarea distintas mediante una mask y un objective diferentes.
